\section*{Bab \ref{ch:g6:symmetry} --- Simetri Sumbu}

\begin{solution}{exo:g6:symmetry:1}
$A'$ berada dua petak di sebelah \emph{kanan} $d$, pada baris yang
sama dengan $A$; $B'$ berada lima petak di sebelah kanan, satu petak
lebih tinggi --- letak yang tercermin.
\end{solution}

\begin{solution}{exo:g6:symmetry:2}
\omterm{def:g6:lines:perp}{Garis tegak lurus} terhadap $d$ melalui $M$, dengan kaki $H$; lalu $M'$
di sisi yang lain dengan $HM' = MH = 3$ cm. Jarak $MM'$ adalah
$3 + 3 = 6$ cm.
\end{solution}

\begin{solution}{exo:g6:symmetry:3}
Sumbu tegak (seperti A): A, H, I, M, O, T, U, V, W, X, Y.
Sumbu mendatar (seperti B): B, C, D, E, H, I, K, O, X.
Keduanya: H, I, O, X.
(Jawaban persisnya sedikit bergantung pada cara hurufnya digambar.)
\end{solution}

\begin{solution}{exo:g6:symmetry:4}
\omterm{def:g6:shapes:triangles}{Segitiga sama sisi}: $3$ sumbu. Persegi panjang: $2$ (\omterm{def:g6:symmetry:bisector}{garis sumbu}
sisi-sisinya). \omterm{def:g6:shapes:quadrilaterals}{Belah ketupat}: $2$ (yaitu diagonalnya). Persegi: $4$.
Lingkaran: tak hingga banyaknya (setiap garis yang melalui pusatnya).
\end{solution}

\begin{solution}{exo:g6:symmetry:5}
\emph{Salah.} Melipat persegi panjang (yang bukan persegi) sepanjang
diagonalnya tidak membuat kedua paruhnya berimpit: kedua segitiganya
punya bentuk yang sama tetapi duduk berbeda --- cobalah dengan
selembar kertas. Hanya pada persegi diagonalnya berlaku sebagai sumbu.
\end{solution}

\begin{solution}{exo:g6:symmetry:6}
Pencerminan mempertahankan panjang, sudut dan \omterm{def:g4:measure:perimeter}{keliling}
(\cref{prop:g6:symmetry:preserved}): $A'B' = 4$ cm,
$\widehat{B'A'C'} = 70^\circ$, \omterm{def:g4:measure:perimeter}{keliling} $A'B'C' = 12$ cm.
\end{solution}

\begin{solution}{exo:g6:symmetry:7}
Lukisan dengan jangka seperti pada \cref{met:g6:symmetry:bisector}.
Untuk sembarang $M$ pada garis itu, penggaris membenarkan $MA = MB$
(\cref{prop:g6:symmetry:bisector}).
\end{solution}

\begin{solution}{exo:g6:symmetry:8}
Semua titik yang berjarak sama dari $A$ dan $B$ membentuk \omterm{def:g6:symmetry:bisector}{garis sumbu}
$[AB]$ (\cref{prop:g6:symmetry:bisector}). Air mancurnya dapat
diletakkan di mana pun pada garis itu (selama masih di dalam kebun).
\end{solution}

\begin{solution}{exo:g6:symmetry:9}
Setiap \omterm{def:g5:solids:def}{titik sudut} dicerminkan \omterm{def:g6:lines:perp}{tegak lurus} terhadap sumbunya: untuk
sumbu $45^\circ$, titik yang berjarak $k$ petak di kanan sumbu
mendarat $k$ petak di atasnya (dan sebaliknya) --- pergeseran mendatar
dan tegaknya bertukar. Cerminkan \omterm{def:g5:solids:def}{titik sudutnya} satu per satu, lalu
hubungkan berurutan.
\end{solution}

\begin{solution}{exo:g6:symmetry:10}
\omterm{def:g6:symmetry:def}{Bayangan} $M$ adalah titik potong kedua lingkaran itu dengan \omterm{def:g6:lines:perp}{garis
tegak lurus} terhadap $d$ melalui $M$ --- ia terletak pada lingkaran.
Alasannya: pencerminan mempertahankan jarak dan membiarkan $O$ (yang
terletak pada $d$) tetap di tempatnya, jadi bayangan setiap titik yang berjarak $r$
dari $O$ tetap berjarak $r$ dari $O$: lingkarannya terpetakan ke
dirinya sendiri.
\end{solution}

\begin{solution}{exo:g6:symmetry:11}
Setiap titik potong $P$ digambar dengan bukaan jangka yang sama
$\rho$ dari $A$ dan dari $B$: jadi $PA = \rho = PB$. Menurut
\cref{prop:g6:symmetry:bisector}, titik yang berjarak sama dari $A$
dan $B$ terletak pada \omterm{def:g6:symmetry:bisector}{garis sumbu} $[AB]$; dua titik semacam itu
menentukan garisnya.
\end{solution}

\begin{solution}{exo:g6:symmetry:12}
Untuk sembarang titik pantul $P$ pada dinding, $PB = PB'$
(pencerminan terhadap $d$ mempertahankan jarak), jadi panjang
jalurnya $AP + PB = AP + PB'$. Garis patah $A \to P \to B'$ paling
pendek ketika ia lurus, yaitu ketika $P$ terletak pada \omterm{def:g6:lines:objects}{ruas garis}
$[AB']$. Jadi titik pantul terbaik adalah perpotongan $[AB']$ dengan dindingnya.
\end{solution}

\begin{solution}{pb:g6:symmetry:1}
\textbf{1.} Lukisannya seperti pada \cref{met:g6:symmetry:bisector};
untuk sembarang $M$ pada garis itu, penggaris membenarkan $MA = MB$
(\cref{prop:g6:symmetry:bisector}).

\textbf{2.} Pertanyaan air mancurnya memerlukan arah yang kedua:
\emph{semua} titik yang berjarak sama terletak pada \omterm{def:g6:symmetry:bisector}{garis sumbunya},
jadi tempat yang mungkin tepatnya garis itu dan tidak ada yang lain.
Dan titik yang tidak pada \omterm{def:g6:symmetry:bisector}{garis sumbunya} karena itu \emph{tidak}
berjarak sama: ia jelas lebih dekat ke salah satu pohonnya.

\textbf{3.} Kedua jalurnya berbagi ruas $[MP]$, dan $PB = PA$ karena
$P$ terletak pada \omterm{def:g6:symmetry:bisector}{garis sumbunya}. Jadi jalur $M \to P \to B$ sama
panjang dengan jalur $M \to P \to A$, yaitu $MP + PA$. Sekarang $MB$
adalah rute lurus ke $B$ --- di sini ia \emph{memang} jalur lewat
$P$, jadi $MB = MP + PA$; dan rute lurus ke $A$ lebih pendek daripada
memutar lewat $P$: $MA < MP + PA = MB$. Setiap titik di sisi $A$ dari
\omterm{def:g6:symmetry:bisector}{garis sumbunya} lebih dekat ke $A$.

\textbf{4.} Batasnya adalah \omterm{def:g6:symmetry:bisector}{garis sumbu} $[AB]$: sekolah $A$ tepat
melayani paruh kota di sisinya sendiri terhadap garis itu (pertanyaan
3), sekolah $B$ paruh yang lain, dan anak yang tinggal tepat pada
garisnya boleh memilih salah satu.

\textbf{5.} $O$ terletak pada \omterm{def:g6:symmetry:bisector}{garis sumbu} $[AB]$, jadi $OA = OB$;
dan pada \omterm{def:g6:symmetry:bisector}{garis sumbu} $[BC]$, jadi $OB = OC$
(\cref{prop:g6:symmetry:bisector}).

\textbf{6.} Dari $OA = OB$ dan $OB = OC$: $OA = OC$. Menurut arah
kedua \cref{prop:g6:symmetry:bisector}, titik yang berjarak sama dari
$A$ dan $C$ terletak pada \omterm{def:g6:symmetry:bisector}{garis sumbu} $[AC]$: \omterm{def:g6:symmetry:bisector}{garis sumbu} ketiga
melalui $O$ tanpa diminta. Ketiga \omterm{def:g6:symmetry:bisector}{garis sumbu} sisi segitiga mana pun
\emph{berpotongan di satu titik}.

\textbf{7.} Lingkaran berpusat $O$ dengan \omterm{def:g3:shapes:shapes}{jari-jari} $OA$ terdiri atas
semua titik yang berjarak $OA$ dari $O$
(\cref{def:g6:lines:circle}); karena $OB = OA$ dan $OC = OA$, titik
$B$ dan $C$ terletak padanya. Satu pusat, satu \omterm{def:g3:shapes:shapes}{jari-jari}, tiga titik
terlayani: itulah lingkaran yang melalui $A$, $B$ dan $C$.

\textbf{8.} Dua \omterm{def:g6:symmetry:bisector}{garis sumbu} sudah cukup karena yang ketiga otomatis
(pertanyaan 6): titik potong keduanya sudah memenuhi ketiga kesamaan
jarak itu. Lingkarannya digambar dengan pusat $O$ dan jangka yang
dibuka dari $O$ ke salah satu \omterm{def:g5:solids:def}{titik sudutnya}.

\textbf{9.} Resepnya: tandailah tiga titik $A$, $B$, $C$ pada busur
yang selamat; gambarlah \omterm{def:g6:symmetry:bisector}{garis sumbu} tali busur $[AB]$ dan $[BC]$;
titik potongnya adalah pusat piring itu, dan jaraknya ke $A$ adalah
\omterm{def:g3:shapes:shapes}{jari-jarinya}. Cara ini berhasil karena tepi piringnya adalah lingkaran
yang melalui ketiga titik yang ditandai --- dan pertanyaan 8 mencari
\emph{satu-satunya} lingkaran yang melalui ketiganya. Busur ujinya terbangun
kembali dengan tepat.

\textbf{10.} Kalau $A$, $B$, $C$ segaris, \omterm{def:g6:symmetry:bisector}{garis sumbu} $[AB]$ dan
\omterm{def:g6:symmetry:bisector}{garis sumbu} $[BC]$ sama-sama \omterm{def:g6:lines:perp}{tegak lurus} terhadap garis yang sama
$(AB) = (BC)$ --- jadi keduanya \omterm{def:g6:lines:perp}{sejajar} satu sama lain
(\cref{prop:g6:lines:facts}) dan, karena berbeda, tidak pernah
berpotongan. Tidak ada titik yang berjarak sama dari ketiganya, dan
tidak ada lingkaran yang melalui tiga titik yang segaris.

\textbf{11.} \omterm{def:g4:geometry:polygon}{Segi lima} beraturan: $5$ sumbu, masing-masing melalui
satu \omterm{def:g5:solids:def}{titik sudut} dan titik tengah sisi di hadapannya (banyaknya
\omterm{def:g3:numbers:evenodd}{ganjil}: setiap sumbunya sejenis). \omterm{def:g4:geometry:polygon}{Segi enam} beraturan: $6$ sumbu ---
tiga melalui \omterm{def:g5:solids:def}{titik sudut} yang berhadapan dan tiga melalui titik tengah
sisi yang berhadapan (banyaknya \omterm{def:g3:numbers:evenodd}{genap}: dua jenis, separuh-separuh).
Polanya: \omterm{def:g4:geometry:polygon}{segi banyak} beraturan dengan $n$ sisi punya $n$ sumbu.

\textbf{12.} Guntingannya menembus kedua lapisan sekaligus, jadi kedua
lapisannya kehilangan potongan yang persis sepadan; membuka lipatannya
adalah pencerminan terhadap garis lipatnya, yang mengirim tepi tiap
lapisan ke tepi lapisan yang lain. Lubangnya adalah bayangan cerminnya
sendiri: garis lipatnya adalah \omterm{def:g6:symmetry:axis}{sumbu simetrinya} --- apa pun yang digunting.

\textbf{13.} Lubang yang sudah dibuka punya sedikitnya $2$ \omterm{def:g6:symmetry:axis}{sumbu
simetri}, yaitu kedua garis lipatnya, dan guntingannya muncul dalam $4$
salinan (setiap kali dibuka gambarnya berlipat dua: $1 \to 2 \to 4$).
Melipat lebih banyak kali dengan gaya kepingan salju melipatgandakan
salinan dan sumbunya lebih jauh lagi.

\textbf{14.} $F''$ punya arah hadap yang \emph{sama} dengan $F$ (dua
pembalikan cermin saling meniadakan), ukuran yang sama dan ketinggian
yang sama --- ia sekadar $F$ yang digeser ke kanan. Pergeserannya
berukuran $8$ petak: tepat \emph{dua kali} jarak antara kedua
sumbunya. Dua cermin yang \omterm{def:g6:lines:perp}{sejajar} bersama-sama sama sekali tidak
mencerminkan: keduanya \emph{menggeser}.

\textbf{15.} $F''$ lagi-lagi berukuran sama, tetapi kini ia tampak
\emph{terbalik}, karena terputar setengah putaran terhadap titik
potong kedua sumbunya: setiap titik $F''$ berseberangan diametral dengan titik
$F$ yang bersesuaian melalui pusat itu. Dua cermin yang \omterm{def:g6:lines:perp}{tegak lurus}
menyusun \emph{setengah putaran} --- yaitu simetri pusat pada
\cref{ch:g7:central}.

\end{solution}
